\section{Related Work}
\label{sec:related_work}

% \subsubsection{Data-Driven Approaches}
% 1. Data-driven concept hierarchy generation from text corpora
% TODO: unclear opening
Automatic concept hierarchy generation from text corpora has been approached in various ways.
%% 1.1. Non-FCA data-driven approaches
Hierarchical text classification assigns documents one or more labels organized in a taxonomy~\cite{yu2022constrained}, implicitly creating a hierarchical structure of documents.
Approaches range from regularization-based classifiers~\cite{gopal2013recursive}, embedding-based document-label matching~\cite{li2019hiercon} to neural sequence models~\cite{risch2020hierarchical,yu2022constrained} and \ac{llm}-assisted taxonomy enrichment~\cite{zhang2025teleclass}.

%% 1.2. FCA-based data-driven approaches
Beyond classification-based approaches, research has explored deriving concept hierarchies directly from textual data using \ac{fca}.
Since \ac{fca}-based hierarchy generation can be viewed as conceptual clustering that yields intensional descriptions of abstract concepts via their intents~\cite{cimiano_learning_2005}, text corpora can be represented as a formal contexts in which documents form the objects and extracted features serve as attributes.
Feature sources include semantic dictionaries~\cite{Priss:2004,old:04,Priss:96}, lexical databases~\cite{Priss:96,priss:98,GalnPez2014DiscoveringNS,val:05}, morphological structures~\cite{sporleder:02,petersen:04}, bag-of-words representations capturing the presence of phrases~\cite{Elzinga:09,cigarran_twitter_fca_2016}, and latent topics derived from topic models~\cite{Aznag2014LeveragingFC,wang_topic_feature_2021,wang_topic_hierarchy_2023,hirth_ordinal_2024,gutekunst2025conceptual}.
% \acl{lda} topic features~\cite{wang_topic_feature_2021,wang_topic_hierarchy_2023}, and other topic model features~\cite{Aznag2014LeveragingFC,gutekunst2025conceptual,hirth_ordinal_2024}. % hirth_ordinal_2024 uses own trained topic model refered to as SSH21, but is something trained using non-negative matrix factorization, gutekunst2025conceptual use top2vec, Aznag2014LeveragingFC: compared LDA, CTM, PLSA (found LDA to be best)


% \subsubsection{Knowledge-Driven Approaches}
% 2. Knowledge-based concept hierarchy generation
% 2.1. Constrained Hierarchical Clustering -> explain MLB constraints
%% clustering with constraints
In constrained clustering, prior knowledge is expressed through constraints rather than derived from data.
% , which require objects to be in the same or different clusters, respectively.
Conventional constrained clustering methods typically operate on a single level of granularity~\cite{wagstaff2001constrained}.
%% limits of pairwise constraints: hierarchical clustering
However, pairwise \textit{must-link} and \textit{cannot-link} constraints become ambiguous in hierarchical clustering settings, where objects may belong to the same cluster at one hierarchy level but not at another~\cite{bade_hierarchical_2014}.
To address this limitation, Bade et al.~\cite{bade2006personalized,bade_creating_2008} introduce \ac{mlb} constraints for \ac{hac}.
These constraints encode relative hierarchical relationships through the order in which objects should be merged during agglomeration~\cite{bade2006personalized,bade_creating_2008,bade_hierarchical_2014}.
Existing approaches incorporate \ac{mlb} constraints by modifying similarity measures during merging~\cite{bade2006personalized}, delaying merges until constraint violations are resolved~\cite{zhao2010hierarchical}, or integrating them into matrix-based approaches~\cite{zheng2011semi}.

% In hierarchical clustering, two main approaches exist: partitional and agglomerative hierarchical clustering.
% While partitional hierarchical clustering build a hierarchy by recursively partitioning the data (i.e., top to bottom), agglomerative hierarchical clustering methods build a hierarchy by iteratively merging clusters (i.e., bottom to top)~\cite{zhao2002evaluation}.

% The \ac{hac} algorithm iteratively merges the two most similar clusters into one. The \ac{ihac} extends the \ac{hac} algorithm by not only considering similarity but also constraint violations~\cite{bade_hierarchical_2014}.

% MacLellan et al.~\cite{maclellan2024trestle} incrementally builds concept hierarchies.

% 2.2. Constrained and Knowledge-Guided FCA
Several works also investigate incorporating constraints or expert knowledge directly into \ac{fca}.
% constrained FCA != we use MLB constraints
Constraints have been used to encode mutually exclusive attributes~\cite{hidayat_constrained_fca_2024}
% we use support, i.e., the number of objects sharing a set of attributes, while belohlavek2009formal define importance directly on attributes, or belohlavek2006formal constrain closure operators, or mao_constrained_fca_2017 define a operator S (!= closure operator) to determine interesting attribute sets
and to prune concept lattices by retaining only concepts whose intents satisfy relevance criteria. These criteria have been defined through attribute properties~\cite{belohlavek2011selecting,belohlavek2009formal,hidayat_constrained_fca_2024}, constraints on the closure operator~\cite{belohlavek2006formal}, or additional operators~\cite{mao_constrained_fca_2017}. 
% Attribute Exploration with contradicting experts: != we do not only have shared implications but also some-what shared (i.e., coherence lattice)
Research has also extended Attribute Exploration to settings with multiple, potentially contradictory or incomplete experts, aiming to derive implications shared across different subsets of experts~\cite{felde2021triadic,felde2022attribute}.

%% Our work
However, prior work has incorporated hierarchical \ac{mlb} constraints mainly into traditional clustering methods rather than hierarchy construction via \ac{fca}. 
To the best of our knowledge, no comparison exists between concept lattices derived from probabilistic topic models and those derived from constraints.